\chapter[Lapisan Data: Objek Data, Model Data, dan Tata Kelola Informasi]{Lapisan
Data: Objek Data, Model Data, dan Tata Kelola Informasi Transfer Dana}
\label{ch:data-architecture}
\setlength{\emergencystretch}{2em}

\section*{Tujuan Pembelajaran}
\begin{itemize}
  \item Menjelaskan elemen informasi ArchiMate lintas tingkat abstraksi---objek
        bisnis (\emph{business object}), objek data (\emph{data object}), dan
        representasi (\emph{representation})---serta relasi realisasi dan akses
        yang menautkannya.
  \item Menyusun model data konseptual dan logis untuk kapabilitas sasaran, serta
        memetakan objek data ke proses bisnis yang mengaksesnya.
  \item Menempatkan analisis data dalam kerangka TOGAF ADM \textbf{Fase C
        (Arsitektur Sistem Informasi)}, khususnya sub-fase \textbf{Arsitektur
        Data}: menyusun arsitektur data \emph{baseline} dan \emph{target} serta
        menganalisis kesenjangannya.
  \item Mengaitkan rancangan data dengan konteks tata kelola dan perlindungan
        data pribadi di Indonesia, lalu mengubahnya menjadi narasi analisis untuk
        makalah.
\end{itemize}

Tujuan pertama dan kedua mendukung CLO 3; tujuan ketiga mendukung CLO 2 dan
CLO 3; tujuan keempat mendukung CLO 6 dan CLO 7 karena mahasiswa menautkan
rancangan data dengan regulasi dan argumen akademik.

\section{Ringkasan Awal Bab}

Bab sebelumnya (Bab~\ref{ch:business-architecture}) memetakan lapisan bisnis:
layanan, proses, peran, dan objek bisnis yang mewujudkan rel transfer. Bab ini
turun satu tingkat ke \textbf{lapisan data}: informasi apa yang mengalir dalam
proses transfer, bagaimana ia disusun sebagai objek data, dan bagaimana data itu
ditelusuri kembali ke objek bisnis serta ke proses yang menggunakannya.

Dalam ArchiMate, data adalah bagian dari lapisan aplikasi: \emph{data object}
menempati struktur pasif lapisan aplikasi, sedangkan pada tingkat bisnis
informasi dinyatakan sebagai \emph{business object}. Dalam TOGAF, keduanya
dikerjakan pada \textbf{Fase C}. Bab ini mengambil sisi \textbf{Data}-nya; sisi
\textbf{Application} dibahas pada Bab~\ref{ch:application-architecture}. Keluaran
bab memperkuat bagian \textbf{Metode} dan \textbf{Analisis} makalah: model data,
katalog entitas data, dan matriks data-ke-proses.

\section{Pendahuluan}

Banyak masalah transformasi digital sebenarnya adalah masalah data: definisi yang
tidak konsisten antarkanal, kepemilikan yang tidak jelas, kualitas yang rendah,
dan kepatuhan privasi yang belum melekat. Sasaran transfer real-time tidak dapat
diwujudkan bila data rekening tersebar, hasil penyaringan tidak tercatat rapi,
atau jejak audit tidak lengkap. Karena itu, sebelum berbicara tentang aplikasi
dan teknologi, arsitektur perlu menegaskan \emph{data apa} yang dikelola dan
\emph{aturan main}-nya.

Bagi makalah publikasi, lapisan data memberi dua hal. Pertama, ia membuat analisis
lebih konkret: alih-alih menyebut ``sistem transfer'', penulis dapat menunjukkan
entitas Instruksi Transfer, Catatan Transaksi, dan Hasil Penyaringan beserta
relasinya. Kedua, ia menautkan rancangan dengan regulasi perlindungan data
pribadi \cite{uupdp2022} dan tata kelola data lintas instansi \cite{satudata},
sehingga rekomendasi tidak berhenti sebagai gambar teknis tetapi menjadi argumen
yang sadar-konteks.

\section{Skenario Pembuka dan Pertanyaan Pemandu}

Setelah proses transfer Bank Wanua dipetakan, tim data menemukan bahwa istilah
``rekening'', ``limit'', dan ``status transaksi'' didefinisikan berbeda di kanal
mobile, internet banking, dan sistem inti. Rekonsiliasi manual sering gagal karena
catatan transaksi tidak memiliki kunci yang seragam. Ketika divisi kepatuhan
meminta jejak audit satu transfer, tim harus menyusun ulang data dari beberapa
sistem. Persoalannya bukan ketiadaan data, melainkan ketiadaan \emph{model data}
yang disepakati.

\subsection*{Pertanyaan Pemandu}
\begin{enumerate}
  \item Objek data apa yang paling menentukan bagi kapabilitas transfer, dan dari
        objek bisnis mana ia diturunkan?
  \item Bagaimana objek data saling berhubungan (rujukan, komposisi, dependensi)?
  \item Proses mana yang menulis dan membaca tiap objek data?
  \item Aturan kualitas, kepemilikan, dan privasi apa yang harus melekat pada data?
  \item Bagaimana kondisi data saat ini berbeda dari kondisi target, dan apa
        kesenjangannya?
\end{enumerate}

\section{Mengapa Bab Ini Penting bagi Pembaca Non-TI}

Pembaca non-TI kerap menyerahkan urusan data kepada tim teknis. Namun keputusan
paling menentukan tentang data bersifat manajerial: siapa pemilik data, sampai
level apa data dibagikan, dan bagaimana risiko privasi dikelola. Model data yang
baik membuat keputusan tersebut terlihat dan dapat didiskusikan tanpa masuk ke
detail basis data. Dengan tiga notasi sederhana---objek bisnis, objek data, dan
representasi---seorang analis dapat menjelaskan bagaimana informasi mengalir dan
di mana risikonya, lalu menautkannya ke argumen makalah.

\section{Konsep Inti dan Glosarium Mini}

\subsection{Objek Data dan Informasi dalam ArchiMate}

ArchiMate memodelkan informasi pada beberapa tingkat abstraksi. \textbf{Objek
bisnis} (\emph{business object}) menyatakan informasi pada tingkat konseptual:
``apa'' yang dikelola organisasi, lepas dari teknologi (misalnya Instruksi
Transfer). \textbf{Objek data} (\emph{data object}) menyatakan informasi pada
tingkat logis/aplikasi: bagaimana data disusun dan diproses aplikasi (misalnya
Data Instruksi Transfer). \textbf{Representasi} (\emph{representation}) adalah
perwujudan yang dapat dipersepsi manusia, seperti struk atau laporan. Objek data
\emph{diturunkan dari} objek bisnis: dalam notasi, objek data
\textbf{merealisasikan} objek bisnis.

\subsection{Notasi Elemen Data di Archi}

Gambar~\ref{fig:data-legend} menampilkan tiga notasi informasi di Archi. Ketiganya
memakai bentuk persegi panjang dengan penanda kepala yang berbeda; objek data
berwarna kebiruan (lapisan aplikasi), sedangkan objek bisnis dan representasi
berwarna kuning (lapisan bisnis).

\begin{figure}[htbp]
\centering
\includegraphics[width=0.85\textwidth]{data-legend.pdf}
\caption{Notasi informasi ArchiMate: objek bisnis (konseptual), objek data
(logis), dan representasi (perwujudan).}
\label{fig:data-legend}
\end{figure}

\begin{table}[htbp]
\centering
\scriptsize
\caption{Elemen informasi ArchiMate, kegunaan, dan contohnya.}
\label{tab:data-elements}
\begin{tabularx}{\textwidth}{@{}p{0.20\textwidth} p{0.20\textwidth} X@{}}
\hline
\textbf{Elemen} & \textbf{Tingkat \& kegunaan} & \textbf{Contoh} \\
\hline
Business Object & Konseptual---``apa'' informasi bisnis, dasar model data
konseptual. & Instruksi Transfer, Rekening Nasabah, Catatan Transaksi, Hasil
Penyaringan; (restoran) Pesanan. \\
Data Object & Logis/aplikasi---bagaimana data disusun \& diproses; mewujudkan
objek bisnis; ditautkan ke proses via objek bisnis. & Data Instruksi Transfer,
Data Rekening, Data Transaksi (Ledger), Data Hasil AML; (restoran) Data Pesanan. \\
Representation & Perwujudan yang dapat dipersepsi (dokumen, laporan, layar). &
Struk/bukti transaksi, laporan rekonsiliasi; (restoran) Struk Pembayaran. \\
\hline
\end{tabularx}
\end{table}

\subsection{Tiga Tingkat Model Data}

Model data biasa dibedakan menjadi tiga tingkat (Gambar~\ref{fig:data-levels}).
\textbf{Konseptual} menjawab entitas informasi utama dan relasinya, dinyatakan
sebagai objek bisnis. \textbf{Logis} merinci atribut dan struktur tanpa terikat
teknologi, dinyatakan sebagai objek data. \textbf{Fisik} adalah wujud tersimpan
pada teknologi (tabel, berkas), dinyatakan sebagai \emph{artifact} pada lapisan
teknologi (Bab~\ref{ch:business-architecture} mempratinjaunya). Bab ini berfokus
pada dua tingkat teratas.

\begin{figure}[htbp]
\centering
\scalebox{0.9}{
\begin{tikzpicture}[
  font=\footnotesize,
  lvl/.style={rectangle, rounded corners=1.5mm, thick, align=center,
    minimum height=11mm, text width=44mm},
  konsep/.style={lvl, draw=praditagreen, fill=praditagreen!14},
  logis/.style={lvl, draw=blue!70!black, fill=blue!12},
  fisik/.style={lvl, draw=gray!60, fill=gray!14},
  link/.style={-Latex, thick, draw=gray!60}
]
  \node[konsep] (k) {\textbf{Konseptual}\\{\scriptsize Business Object}};
  \node[logis, below=8mm of k] (l) {\textbf{Logis}\\{\scriptsize Data Object}};
  \node[fisik, below=8mm of l] (f) {\textbf{Fisik}\\{\scriptsize Artifact (Bab~7)}};
  \draw[link] (l) -- node[right]{\scriptsize merealisasikan} (k);
  \draw[link] (f) -- node[right]{\scriptsize merealisasikan} (l);
  \node[align=left, text width=52mm, right=8mm of k] {Entitas informasi \& relasi utama};
  \node[align=left, text width=52mm, right=8mm of l] {Atribut, struktur, kunci---lepas teknologi};
  \node[align=left, text width=52mm, right=8mm of f] {Tabel, berkas, skema pada teknologi};
\end{tikzpicture}
}
\caption{Tiga tingkat model data dan padanannya di ArchiMate.}
\label{fig:data-levels}
\end{figure}

\subsection{TOGAF ADM Fase C: Arsitektur Data}

Dalam siklus TOGAF ADM, data dikerjakan pada \textbf{Fase C (Arsitektur Sistem
Informasi)}, yang terbagi menjadi \emph{Arsitektur Data} dan \emph{Arsitektur
Aplikasi}. Bab ini mengambil sisi data. Seperti Fase B, Fase C menghasilkan dua
potret---\emph{baseline} dan \emph{target}---lalu menganalisis kesenjangannya
\cite{togaf10}. Keluaran utamanya adalah katalog entitas data, diagram model
data, dan matriks entitas-data terhadap fungsi/proses bisnis.
Tabel~\ref{tab:adm-phase-c-data} memetakan langkahnya ke kasus Bank Wanua.

\begin{table}[htbp]
\centering
\scriptsize
\caption{Pemetaan TOGAF ADM Fase C (Arsitektur Data) ke artefak Bank Wanua.}
\label{tab:adm-phase-c-data}
\begin{tabularx}{\textwidth}{@{}p{0.26\textwidth} X@{}}
\hline
\textbf{Langkah Fase C (Data)} & \textbf{Artefak / keluaran pada kasus Bank Wanua} \\
\hline
Ruang lingkup dan prinsip data & Lingkup data transfer; prinsip data sebagai
aset, minimisasi data pribadi, dan jejak audit melekat. \\
Arsitektur data baseline & Definisi rekening/limit/status berbeda antarkanal;
catatan transaksi tanpa kunci seragam; jejak audit tersebar. \\
Arsitektur data target & Objek data logis yang mewujudkan objek bisnis; kunci
transaksi seragam; Data Log Audit dan Data Hasil AML terstruktur. \\
Analisis kesenjangan & Selisih baseline--target: entitas yang perlu
distandardisasi, disatukan, atau dibangun kualitasnya. \\
Masukan ke fase berikutnya & Kebutuhan lapisan aplikasi: layanan data, komponen
integrasi, dan mesin rekonsiliasi (Bab~\ref{ch:application-architecture}). \\
\hline
\end{tabularx}
\end{table}

Selisih baseline--target divisualkan pada Gambar~\ref{fig:data-asis} (as-is) dan
Gambar~\ref{fig:data-tobe} (to-be): dari data yang terfragmentasi dan belum
selaras, menuju objek data yang mewujudkan objek bisnis, berkunci transaksi
seragam, dan berjejak audit melekat.

\begin{figure}[htbp]
\centering
\includegraphics[width=\textwidth,height=0.4\textheight,keepaspectratio]{data-asis.pdf}
\caption{Arsitektur data \emph{as-is} (baseline): definisi berbeda antarkanal,
catatan transaksi tanpa kunci seragam, jejak audit tersebar, hasil AML manual.}
\label{fig:data-asis}
\end{figure}

\begin{figure}[htbp]
\centering
\includegraphics[width=\textwidth,height=0.4\textheight,keepaspectratio]{data-tobe.pdf}
\caption{Arsitektur data \emph{to-be} (target): objek data \emph{mewujudkan}
objek bisnis (definisi selaras), kunci transaksi seragam, dan jejak audit serta
hasil AML terstruktur.}
\label{fig:data-tobe}
\end{figure}

\subsection{Contoh Interaksi Antar-Elemen Data}

Sebelum masuk ke kasus Bank Wanua, Gambar~\ref{fig:data-contoh} memakai contoh
sederhana restoran cepat saji untuk menunjukkan cara ketiga notasi berinteraksi:
objek data \emph{Data Pesanan} merealisasikan objek bisnis \emph{Pesanan}---yang
diakses proses \emph{Menyajikan Pesanan}---dan objek bisnis itu diwujudkan bagi
pelanggan sebagai representasi \emph{Struk Pembayaran}, sementara data pendukung
(menu, pelanggan, pembayaran) melengkapinya.

\begin{figure}[htbp]
\centering
\includegraphics[width=\textwidth,height=0.42\textheight,keepaspectratio]{data-contoh.pdf}
\caption{Contoh interaksi antar-elemen data pada kasus restoran cepat saji.}
\label{fig:data-contoh}
\end{figure}

\subsection*{Glosarium Mini}
\begin{description}
  \item[Objek bisnis (business object)] Informasi tingkat konseptual yang dikelola
        organisasi, misalnya Instruksi Transfer.
  \item[Objek data (data object)] Representasi logis data yang diproses aplikasi
        dan merealisasikan objek bisnis, misalnya Data Instruksi Transfer.
  \item[Representasi (representation)] Perwujudan informasi yang dapat dipersepsi,
        misalnya Struk Pembayaran.
  \item[Relasi akses] Relasi yang menyatakan proses membaca atau menulis objek
        data.
  \item[Model data konseptual/logis] Peta entitas informasi dan relasinya
        (konseptual) atau struktur atribut dan kuncinya (logis).
  \item[Tata kelola data] Aturan kepemilikan, kualitas, definisi, dan akses data.
\end{description}

\section{Kerangka Berpikir}

Alur berpikir bab ini bergerak dari objek bisnis menuju arsitektur data
(Gambar~\ref{fig:data-flow}). Objek bisnis dari lapisan bisnis menjadi titik awal;
tiap objek bisnis diturunkan menjadi objek data logis; objek data dihubungkan
menjadi model data; model data dipetakan ke proses melalui relasi akses; dan
aturan tata kelola melekatkan kualitas serta privasi. Hasilnya menjadi bahan
analisis dan draf metode makalah.

\begin{figure}[htbp]
\centering
\scalebox{0.86}{
\begin{tikzpicture}[
  box/.style={rectangle, rounded corners=1.5mm, draw=blue!70!black,
    fill=blue!8, thick, align=center, text width=30mm, minimum height=11mm},
  outnode/.style={rectangle, rounded corners=1.5mm, draw=orange!80!black,
    fill=orange!16, thick, align=center, text width=32mm, minimum height=11mm},
  link/.style={-Latex, thick, draw=blue!70!black}
]
  \node[box] (bo) {Objek bisnis\\(dari Bab 5)};
  \node[box, right=7mm of bo] (do) {Objek data\\logis};
  \node[box, right=7mm of do] (model) {Model data\\\& relasi};
  \node[box, below=8mm of model] (akses) {Pemetaan\\data-ke-proses};
  \node[outnode, left=7mm of akses] (gov) {Tata kelola\\\& privasi};
  \node[outnode, left=7mm of gov] (an) {Analisis \&\\draf metode};
  \draw[link] (bo) -- (do); \draw[link] (do) -- (model); \draw[link] (model) -- (akses);
  \draw[link] (akses) -- (gov); \draw[link] (gov) -- (an);
  \draw[link, dashed] (an) -- (bo);
\end{tikzpicture}
}
\caption{Alur berpikir lapisan data untuk makalah publikasi.}
\label{fig:data-flow}
\end{figure}

\section{Konteks Indonesia}

Arsitektur data di Indonesia bergerak di antara ambisi berbagi data dan kewajiban
melindunginya. Undang-Undang Pelindungan Data Pribadi menuntut dasar pemrosesan,
minimisasi, pembatasan akses, dan pencatatan penggunaan data pribadi
\cite{uupdp2022}. Penyelenggara sistem elektronik wajib mengelola sistem dan data
secara andal serta menyimpan jejak \cite{pp712019}. Untuk layanan transfer,
standar Open API pembayaran nasional turut membakukan format data pertukaran
\cite{snapbi}. Pada konteks lintas instansi pemerintah, prinsip Satu Data
Indonesia---data yang akurat, mutakhir, terpadu, dan dapat dibagipakaikan dengan
standar dan metadata yang jelas---menjadi acuan tata kelola \cite{satudata}.

Implikasinya bagi model data cukup langsung: entitas yang memuat data pribadi
(Data Nasabah/KYC) memerlukan pembatasan akses dan retensi; jejak audit (Data Log
Audit) menjadi kebutuhan wajib, bukan tambahan; dan definisi data transaksi perlu
seragam agar interoperabel lintas kanal dan mitra.

\section{Studi Kasus Utama: Arsitektur Data Transfer Bank Wanua}

Objek-objek bisnis pada Bab~\ref{ch:business-architecture}---Instruksi Transfer,
Rekening Nasabah, Catatan Transaksi, Hasil Penyaringan, Notifikasi Status, dan
Catatan Rekonsiliasi---diturunkan menjadi objek data logis dan disatukan menjadi
model data pada Gambar~\ref{fig:data-transfer}.

\begin{figure}[htbp]
\centering
\includegraphics[width=\textwidth,height=0.6\textheight,keepaspectratio]{data-transfer.pdf}
\caption{Model arsitektur data (ArchiMate) API Transfer Dana Bank Wanua: objek
data logis (kiri/tengah) yang merealisasikan objek bisnis konseptual (kanan),
beserta relasi antar-data dan data pendukung/master.}
\label{fig:data-transfer}
\end{figure}

\subsection{Cara Membaca Diagram Data}

Gambar~\ref{fig:data-transfer} dibaca dari kanan ke kiri. Kolom kanan (kuning)
memuat \textbf{objek bisnis} konseptual---``apa'' informasi yang dikelola. Kolom
tengah (biru) memuat \textbf{objek data} logis yang \emph{merealisasikan} objek
bisnis di sebelahnya: Data Instruksi Transfer mewujudkan Instruksi Transfer, Data
Transaksi (Ledger) mewujudkan Catatan Transaksi, dan seterusnya. Inilah wujud
prinsip ``data diturunkan dari objek bisnis''. Kolom kiri memuat \textbf{data
pendukung dan master} yang tidak selalu punya padanan objek bisnis langsung: Data
Nasabah (KYC), Referensi Bank (BI-FAST), Data Limit \& Kebijakan, dan Data Log
Audit.

Relasi antar-objek data membentuk model data logis. \emph{Data Transaksi} berisi
(agregasi) \emph{Data Instruksi Transfer}; instruksi merujuk \emph{Data Rekening}
dan dibuat atas nama \emph{Data Nasabah}; sebuah rekening adalah milik seorang
nasabah; \emph{Data Hasil AML} menyaring instruksi dengan menerapkan \emph{Data
Limit \& Kebijakan}; transaksi memicu \emph{Data Notifikasi}; \emph{Data
Rekonsiliasi} mencocokkan transaksi; dan \emph{Data Log Audit} mencatat transaksi.
Label pada tiap garis menyatakan makna relasi sehingga model dapat dibaca sebagai
kalimat.

\subsection{Keterlacakan Data ke Proses}

Model data menjadi lebih berguna ketika ditautkan dengan proses yang memakainya.
Dalam ArchiMate, proses bisnis \emph{mengakses objek bisnis}---bukan objek data
secara langsung---sedangkan objek data \emph{merealisasikan} objek bisnis itu.
Karena itu keterlacakan data-ke-proses ditelusuri berlapis
(Gambar~\ref{fig:data-akses}): proses mengakses (baca/tulis) objek bisnis di kolom
tengah, dan objek data di kolom kanan merealisasikan objek bisnis tersebut.
Misalnya, Inisiasi menulis objek bisnis Instruksi Transfer yang direalisasikan
oleh Data Instruksi Transfer; Validasi \& Otorisasi membaca Rekening Nasabah;
Penyaringan AML \& Limit menulis Hasil Penyaringan; Eksekusi menulis Catatan
Transaksi; dan Rekonsiliasi menulis Catatan Rekonsiliasi. Dengan menelusuri rantai
\emph{akses}~+~\emph{realisasi}, kita mengetahui objek data mana yang dipakai
proses mana tanpa menautkan proses langsung ke objek data.

\begin{figure}[htbp]
\centering
\includegraphics[width=\textwidth,height=0.6\textheight,keepaspectratio]{data-akses.pdf}
\caption{Keterlacakan data-ke-proses berlapis: proses bisnis \emph{mengakses}
(baca/tulis) objek bisnis, dan objek data \emph{merealisasikan} objek bisnis.
Proses tidak mengakses objek data secara langsung.}
\label{fig:data-akses}
\end{figure}

Keterlacakan ini adalah artefak Fase C yang penting. Ia memperlihatkan entitas
paling ``sibuk'' (Catatan Transaksi dan Instruksi Transfer), menandai di mana
kualitas data paling kritis, dan menjadi dasar keputusan kepemilikan data serta
kebutuhan integrasi. Ketika lapisan aplikasi ditambahkan (bab berikutnya), fungsi
aplikasilah yang mengakses objek data secara langsung, melengkapi rantai
proses--layanan--fungsi--data.

\section{Langkah Analisis dan Artefak Pendukung}

\subsection{Langkah 1: Ambil Objek Bisnis dari Lapisan Bisnis}
Mulai dari objek bisnis yang sudah dipetakan di Bab~\ref{ch:business-architecture}.
Objek inilah kandidat entitas informasi utama.

\subsection{Langkah 2: Turunkan Objek Data Logis}
Untuk tiap objek bisnis, tentukan objek data yang merealisasikannya. Tambahkan
data pendukung dan master (misalnya Data Nasabah, Referensi Bank) yang diperlukan.

\subsection{Langkah 3: Susun Model Data}
Hubungkan objek data dengan relasi bermakna (rujukan, komposisi/agregasi,
dependensi). Beri label agar model dapat dibaca sebagai kalimat.

\subsection{Langkah 4: Petakan Data ke Proses}
Nyatakan relasi akses baca/tulis antara proses dan objek data. Tabel matriks
data-ke-proses membantu menemukan entitas paling kritis.

\subsection{Langkah 5: Tetapkan Aturan Tata Kelola}
Untuk tiap entitas, tuliskan pemilik, definisi baku, aturan kualitas, klasifikasi
sensitivitas, retensi, dan dasar pemrosesan bila memuat data pribadi.

\subsection{Langkah 6: Nilai Baseline dan Target}
Bandingkan definisi dan kualitas data saat ini dengan kondisi target, lalu catat
kesenjangan sebagai bahan roadmap.

\subsection{Langkah 7: Kaitkan dengan Regulasi}
Tautkan entitas data pribadi dan jejak audit dengan kewajiban UU PDP dan PP 71,
serta format pertukaran dengan SNAP BI.

\subsection{Langkah 8: Tulis Narasi Analisis Data}
Ubah model dan matriks menjadi paragraf yang menjelaskan keputusan data dan
implikasinya bagi lapisan aplikasi.

\begin{table}[htbp]
\centering
\scriptsize
\caption{Artefak pendukung Bab~6 dan fungsinya dalam makalah.}
\label{tab:artefak-bab5-data}
\begin{tabularx}{\textwidth}{@{}p{0.26\textwidth} p{0.30\textwidth} X@{}}
\hline
\textbf{Artefak} & \textbf{Isi minimum} & \textbf{Dipakai untuk bagian makalah} \\
\hline
Katalog entitas data & Daftar objek bisnis/data, pemilik, definisi & Metode dan
analisis domain data. \\
Diagram model data & Objek data \& relasinya (konseptual/logis) & Analisis dan
lampiran. \\
Matriks data-ke-proses & Akses baca/tulis per proses & Analisis kesenjangan dan
prioritas. \\
Tabel tata kelola \& privasi & Sensitivitas, retensi, dasar pemrosesan & Diskusi,
implikasi, dan konteks Indonesia. \\
\hline
\end{tabularx}
\end{table}

\section{Integrasi ke Makalah Publikasi}

Lapisan data memperkuat bagian \textbf{Metode} dan \textbf{Analisis}. Model data
dan matriks data-ke-proses menjadi bukti bahwa rancangan transfer berpijak pada
informasi yang terdefinisi, bukan asumsi. Contoh paragraf:

\begin{quote}
Analisis data menunjukkan bahwa kapabilitas transfer real-time bertumpu pada enam
entitas inti yang saling terkait, dengan Data Transaksi sebagai simpul terpadat.
Ketiadaan kunci transaksi yang seragam pada kondisi baseline menjelaskan
kegagalan rekonsiliasi otomatis; arsitektur data target menstandardisasi kunci
tersebut dan melekatkan jejak audit sesuai kewajiban perlindungan data.
\end{quote}

\section{Diskusi Kritis, Trade-off, dan Perangkap Umum}

\begin{table}[htbp]
\centering
\scriptsize
\caption{Trade-off dalam arsitektur data.}
\label{tab:tradeoff-bab5-data}
\begin{tabularx}{\textwidth}{@{}p{0.24\textwidth} X X@{}}
\hline
\textbf{Keputusan} & \textbf{Jika terlalu ringan} & \textbf{Jika terlalu berat} \\
\hline
Detail model data & Entitas kabur, relasi tak jelas & Model serumit skema basis
data, sulit dibaca non-TI. \\
Sentralisasi data & Silo data bertahan, definisi berbeda & Titik tunggal
kegagalan dan beban integrasi besar. \\
Data pribadi & Risiko pelanggaran privasi & Fitur macet karena kontrol
berlebihan. \\
Jejak audit & Kepatuhan lemah, sulit ditelusuri & Biaya penyimpanan dan
kompleksitas naik. \\
\hline
\end{tabularx}
\end{table}

Perangkap umum: (1) menyamakan model data dengan skema basis data sehingga terlalu
teknis; (2) melupakan kepemilikan dan definisi baku sehingga silo bertahan; (3)
menambahkan data pribadi tanpa dasar pemrosesan; dan (4) menggambar entitas tanpa
menautkannya ke proses yang memakainya.

\section{Latihan Terapan Individual}

\subsection*{Latihan~6.1: Katalog dan Model Data}
Untuk organisasi studi kasus Anda, turunkan objek data dari objek bisnis, susun
model data logis dengan relasi berlabel, dan tuliskan pemilik serta definisi tiap
entitas.

\subsection*{Latihan~6.2: Matriks Data-ke-Proses}
Buat matriks akses baca/tulis antara proses dan objek data. Tandai entitas paling
kritis dan jelaskan implikasinya bagi kualitas dan integrasi.

\subsection*{Latihan~6.3: Tata Kelola dan Privasi}
Klasifikasikan sensitivitas tiap entitas, tetapkan retensi dan dasar pemrosesan
data pribadi, lalu kaitkan dengan UU PDP dan PP 71.

\subsection*{Latihan~6.4: Baseline, Target, dan Paragraf Analisis}
Bandingkan kondisi data saat ini dan target, catat kesenjangan, lalu tulis satu
paragraf analisis data untuk makalah.

\begin{table}[htbp]
\centering
\scriptsize
\caption{Rubrik latihan Bab~6.}
\label{tab:rubrik-bab5-data}
\begin{tabularx}{\textwidth}{@{}p{0.25\textwidth} X X@{}}
\hline
\textbf{Komponen} & \textbf{Indikator baik} & \textbf{Perlu diperbaiki bila} \\
\hline
Model data & Entitas \& relasi jelas dan berlabel & Sekadar daftar tabel tanpa
relasi. \\
Matriks data-ke-proses & Akses baca/tulis lengkap \& konsisten & Tidak menautkan
data dengan proses. \\
Tata kelola \& privasi & Sensitivitas, retensi, dasar pemrosesan jelas & Data
pribadi tanpa kontrol. \\
Narasi analisis & Menautkan data ke kesenjangan \& keputusan & Deskriptif tanpa
argumen. \\
\hline
\end{tabularx}
\end{table}

\section{Ringkasan, Refleksi, dan Jembatan}

\subsection*{Poin Kunci}
\begin{itemize}
  \item Informasi dimodelkan pada tiga tingkat: objek bisnis (konseptual), objek
        data (logis), dan representasi (perwujudan).
  \item Objek data \emph{diturunkan dari} objek bisnis melalui relasi realisasi.
  \item Model data menjadi berguna ketika dipetakan ke proses lewat relasi akses.
  \item Tata kelola dan privasi (UU PDP, PP 71, SNAP BI, Satu Data) melekat pada
        entitas, bukan tempelan.
  \item Fase C (Data) menghasilkan baseline, target, dan analisis kesenjangan.
\end{itemize}

\subsection*{Uji Mandiri}
\begin{enumerate}
  \item Jelaskan perbedaan objek bisnis, objek data, dan representasi beserta
        contohnya.
  \item Mengapa objek data dikatakan diturunkan dari objek bisnis?
  \item Ambil satu proses pada kasus Anda dan sebutkan objek data yang ditulis dan
        dibacanya.
\end{enumerate}

\subsection*{Jembatan ke Bab Berikutnya}
Bab berikutnya melengkapi Fase C dengan \textbf{lapisan aplikasi}: komponen,
layanan, antarmuka, dan fungsi aplikasi yang \emph{menjalankan} proses bisnis dan
\emph{mengakses} objek data yang telah dimodelkan di sini. Katalog entitas dan
matriks data-ke-proses menjadi masukan langsung untuk memilih komponen aplikasi
dan pola integrasinya.

\section{Bacaan Lanjutan dan Lampiran}

\subsection*{Bacaan Inti}
\begin{itemize}
  \item The Open Group untuk Fase C ADM dan artefak arsitektur data \cite{togaf10}.
  \item Ross, Weill, dan Robertson untuk hubungan integrasi data dengan operating
        model \cite{ross2006enterprise}.
\end{itemize}

\subsection*{Bacaan untuk Memperdalam}
\begin{itemize}
  \item Untuk kepemilikan data pada arsitektur terdistribusi, baca pembahasan
        modularisasi dan kepemilikan data \cite{newman2021building}.
  \item Untuk konteks Indonesia: perlindungan data pribadi \cite{uupdp2022},
        penyelenggaraan sistem elektronik \cite{pp712019}, standar API pembayaran
        \cite{snapbi}, dan Satu Data Indonesia \cite{satudata}.
\end{itemize}

\subsection*{Lampiran Kerja Bab~6}
\begin{description}
  \item[Katalog entitas data] Daftar objek bisnis/data, pemilik, dan definisi.
  \item[Diagram model data] Objek data dan relasinya.
  \item[Matriks data-ke-proses] Akses baca/tulis per proses.
  \item[Tabel tata kelola] Sensitivitas, retensi, dan dasar pemrosesan.
\end{description}
\setlength{\emergencystretch}{0pt}
